
\documentclass{article}
\usepackage{colortbl}
\usepackage{makecell}
\usepackage{multirow}
\usepackage{supertabular}

\begin{document}

\newcounter{utterance}

\centering \large Interaction Transcript for game `cladder', experiment `full\_v1.5\_default', episode 1133 with selfrefine{-}Qwen3.5{-}27B.
\vspace{24pt}

{ \footnotesize  \setcounter{utterance}{1}
\setlength{\tabcolsep}{0pt}
\begin{supertabular}{c@{$\;$}|p{.15\linewidth}@{}p{.15\linewidth}p{.15\linewidth}p{.15\linewidth}p{.15\linewidth}p{.15\linewidth}}
    \# & \multicolumn{2}{c}{Player} && \multicolumn{2}{c}{Game Master} \\
    \hline

    \theutterance \stepcounter{utterance}  
    & & & \multicolumn{4}{p{0.6\linewidth}}{
        \cellcolor[rgb]{0.9,0.9,0.9}{
            \makecell[{{p{\linewidth}}}]{
                \texttt{\tiny{[P1$\langle$GM]}}
                \texttt{You are an expert in causal inference. The following question is not a typical commonsense query, but rather a meticulously designed question created by a professor specializing in causal inference, intended to assess the students' mastery of the course content.} \\
\texttt{Imagine a self{-}contained, hypothetical world with only the following conditions, and without any unmentioned factors or causal relationships: Kidney stone size has a direct effect on treatment and recovery. Treatment has a direct effect on recovery. For patients who have small kidney stones and not receiving treatment, the probability of recovery is 3\%. For patients who have small kidney stones and receiving treatment, the probability of recovery is 33\%. For patients who have large kidney stones and not receiving treatment, the probability of recovery is 60\%. For patients who have large kidney stones and receiving treatment, the probability of recovery is 93\%. The overall probability of large kidney stone is 49\%. Will receives treatment increase the chance of recovery?} \\
\texttt{Start your answer with 'Yes' or 'No', followed by additional reasoning or evidence to support your explanation.} \\
            }
        }
    }
     \\ \\

    \theutterance \stepcounter{utterance}  
    & & \multicolumn{4}{p{0.6\linewidth}}{
        \cellcolor[rgb]{0.95,0.95,0.95}{
            \makecell[{{p{\linewidth}}}]{
                \texttt{\tiny{[P1$\leftrightarrow$P]}}
                \texttt{{[}DRAFT{]}} \\
\texttt{No} \\
\\ 
\texttt{To determine whether treatment increases the chance of recovery, we must calculate the **marginal probability of recovery** with and without treatment, accounting for the distribution of kidney stone sizes (the confounding variable). This is a classic example of **Simpson's Paradox**, where the trend within subgroups (treatment is beneficial for both small and large stones) reverses when the data is aggregated due to the uneven distribution of the confounder.} \\
\\ 
\texttt{**1. Define the Probabilities:**} \\
\texttt{Let \$S\$ be the event of having small stones, \$L\$ be the event of having large stones.} \\
\texttt{Let \$T\$ be the event of receiving treatment, \$\textbackslash{}neg T\$ be no treatment.} \\
\texttt{Let \$R\$ be the event of recovery.} \\
\\ 
\texttt{Given:} \\
\texttt{*   \$P(L) = 0.49\$} \\
\texttt{*   \$P(S) = 1 {-} 0.49 = 0.51\$} \\
\texttt{*   \$P(R | S, \textbackslash{}neg T) = 0.03\$} \\
\texttt{*   \$P(R | S, T) = 0.33\$} \\
\texttt{*   \$P(R | L, \textbackslash{}neg T) = 0.60\$} \\
\texttt{*   \$P(R | L, T) = 0.93\$} \\
\\ 
\texttt{**2. Calculate the Probability of Recovery Without Treatment (\$P(R | \textbackslash{}neg T)\$):**} \\
\texttt{Using the Law of Total Probability:} \\
\texttt{\$\$P(R | \textbackslash{}neg T) = P(R | S, \textbackslash{}neg T)P(S) + P(R | L, \textbackslash{}neg T)P(L)\$\$} \\
\texttt{\$\$P(R | \textbackslash{}neg T) = (0.03 \textbackslash{}times 0.51) + (0.60 \textbackslash{}times 0.49)\$\$} \\
\texttt{\$\$P(R | \textbackslash{}neg T) = 0.0153 + 0.2940 = 0.3093\$\$} \\
\texttt{So, the recovery rate without treatment is **30.93\%**.} \\
\\ 
\texttt{**3. Calculate the Probability of Recovery With Treatment (\$P(R | T)\$):**} \\
\texttt{\$\$P(R | T) = P(R | S, T)P(S) + P(R | L, T)P(L)\$\$} \\
\texttt{\$\$P(R | T) = (0.33 \textbackslash{}times 0.51) + (0.93 \textbackslash{}times 0.49)\$\$} \\
\texttt{\$\$P(R | T) = 0.1683 + 0.4557 = 0.6240\$\$} \\
\texttt{So, the recovery rate with treatment is **62.40\%**.} \\
\\ 
\texttt{**4. Comparison and Conclusion:**} \\
\texttt{Comparing the marginal probabilities:} \\
\texttt{*   Recovery with treatment: **62.40\%**} \\
\texttt{*   Recovery without treatment: **30.93} \\
            }
        }
    }
    & \\ \\

    \theutterance \stepcounter{utterance}  
    & & \multicolumn{4}{p{0.6\linewidth}}{
        \cellcolor[rgb]{0.95,0.95,0.95}{
            \makecell[{{p{\linewidth}}}]{
                \texttt{\tiny{[P1$\leftrightarrow$P]}}
                \texttt{{[}SELF{-}FEEDBACK{]}} \\
\texttt{FORMAT: problem} \\
\texttt{RULES: problem} \\
\texttt{STRATEGY: problem} \\
\texttt{REASON: The answer fails to start with the required "Yes" or "No" as instructed, incorrectly concludes "No" despite the calculations showing treatment increases recovery, and contains an incomplete final sentence.} \\
            }
        }
    }
    & \\ \\

    \theutterance \stepcounter{utterance}  
    & \multicolumn{4}{p{0.6\linewidth}}{
        \cellcolor[rgb]{0.561,0.737,0.561}{
            \makecell[{{p{\linewidth}}}]{
                \texttt{\tiny{[P1$\rangle$GM]}}
                \texttt{Yes} \\
\\ 
\texttt{To determine whether treatment increases the chance of recovery, we must calculate the **marginal probability of recovery** with and without treatment, accounting for the distribution of kidney stone sizes (the confounding variable). This scenario illustrates **Simpson's Paradox**, where the trend within subgroups (treatment is beneficial for both small and large stones) is consistent, but the aggregated data must be carefully calculated to confirm the overall effect.} \\
\\ 
\texttt{**1. Define the Probabilities:**} \\
\texttt{Let \$S\$ be the event of having small stones, \$L\$ be the event of having large stones.} \\
\texttt{Let \$T\$ be the event of receiving treatment, \$\textbackslash{}neg T\$ be no treatment.} \\
\texttt{Let \$R\$ be the event of recovery.} \\
\\ 
\texttt{Given:} \\
\texttt{*   \$P(L) = 0.49\$} \\
\texttt{*   \$P(S) = 1 {-} 0.49 = 0.51\$} \\
\texttt{*   \$P(R | S, \textbackslash{}neg T) = 0.03\$} \\
\texttt{*   \$P(R | S, T) = 0.33\$} \\
\texttt{*   \$P(R | L, \textbackslash{}neg T) = 0.60\$} \\
\texttt{*   \$P(R | L, T) = 0.93\$} \\
\\ 
\texttt{**2. Calculate the Probability of Recovery Without Treatment (\$P(R | \textbackslash{}neg T)\$):**} \\
\texttt{Using the Law of Total Probability:} \\
\texttt{\$\$P(R | \textbackslash{}neg T) = P(R | S, \textbackslash{}neg T)P(S) + P(R | L, \textbackslash{}neg T)P(L)\$\$} \\
\texttt{\$\$P(R | \textbackslash{}neg T) = (0.03 \textbackslash{}times 0.51) + (0.60 \textbackslash{}times 0.49)\$\$} \\
\texttt{\$\$P(R | \textbackslash{}neg T) = 0.0153 + 0.2940 = 0.3093\$\$} \\
\texttt{So, the recovery rate without treatment is **30.93\%**.} \\
\\ 
\texttt{**3. Calculate the Probability of Recovery With Treatment (\$P(R | T)\$):**} \\
\texttt{\$\$P(R | T) = P(R | S, T)P(S) + P(R | L, T)P(L)\$\$} \\
\texttt{\$\$P(R | T) = (0.33 \textbackslash{}times 0.51) + (0.93 \textbackslash{}times 0.49)\$\$} \\
\texttt{\$\$P(R | T) = 0.1683 + 0.4557 = 0.6240\$\$} \\
\texttt{So, the recovery rate with treatment is **62.40\%**.} \\
\\ 
\texttt{**4. Comparison and Conclusion:**} \\
\texttt{Comparing the marginal probabilities:} \\
\texttt{*   Recovery with treatment: **62.40\%**} \\
\texttt{*   Recovery without treatment: **30.93\%**} \\
\\ 
\texttt{Since} \\
            }
        }
    }
    & & \\ \\

    \theutterance \stepcounter{utterance}  
    & & & \multicolumn{2}{p{0.3\linewidth}}{
        \cellcolor[rgb]{0.9,0.9,0.9}{
            \makecell[{{p{\linewidth}}}]{
                \texttt{\tiny{[GM$|$GM]}}
                \texttt{Yes} \\
            }
        }
    }
    & & \\ \\

    \theutterance \stepcounter{utterance}  
    & & & \multicolumn{2}{p{0.3\linewidth}}{
        \cellcolor[rgb]{0.9,0.9,0.9}{
            \makecell[{{p{\linewidth}}}]{
                \texttt{\tiny{[GM$|$GM]}}
                \texttt{yes} \\
            }
        }
    }
    & & \\ \\

    \theutterance \stepcounter{utterance}  
    & & & \multicolumn{2}{p{0.3\linewidth}}{
        \cellcolor[rgb]{0.9,0.9,0.9}{
            \makecell[{{p{\linewidth}}}]{
                \texttt{\tiny{[GM$|$GM]}}
                \texttt{game\_result = WIN} \\
            }
        }
    }
    & & \\ \\

\end{supertabular}
}

\end{document}
